\section*{ЗАКЛЮЧЕНИЕ}
\addcontentsline{toc}{section}{ЗАКЛЮЧЕНИЕ}

В результате выполнения выпускной квалификационной работы была разработана платформа онлайн-бронирования услуг Easyappoint. Система предназначена для автоматизации записи клиентов к частному мастеру и позволяет заменить ручное согласование времени через звонки и сообщения централизованным web-интерфейсом.

В ходе работы были получены следующие результаты:
\begin{enumerate}
\item Проведен анализ предметной области онлайн-бронирования услуг. Выявлены основные проблемы ручной записи: риск двойного бронирования, отсутствие единого источника данных, зависимость от постоянного участия мастера и сложность учета внешнего календаря.
\item Сформулированы требования к программной системе. Определены роли мастера, клиента и администратора, требования к рабочей доступности, приглашениям, бронированиям, уведомлениям, надежности, безопасности и пользовательскому интерфейсу.
\item Спроектирована архитектура платформы. Выделены клиентское web-приложение, серверный API-сервис, сервис Telegram-уведомлений, база данных PostgreSQL, интеграция с Google Calendar и событийный обмен через Kafka.
\item Спроектирована и реализована модель данных. В базе данных предусмотрены таблицы пользователей, ролей, сессий, мастеров, клиентов, связей мастер-клиент, приглашений, расписания, бронирований, подключений календаря, уведомлений и outbox-событий.
\item Реализована серверная часть платформы на Kotlin с использованием Ktor. Серверная часть обрабатывает регистрацию, аутентификацию, управление расписанием, приглашениями, бронированиями, внешним календарем и каналами уведомлений.
\item Реализовано клиентское web-приложение на React и TypeScript. Интерфейс содержит ролевые кабинеты мастера, клиента и администратора, экраны работы с расписанием, клиентами, бронированиями, профилем, сервисами и уведомлениями.
\item Реализован сервис Telegram-уведомлений. Он получает события из Kafka, определяет тип события и отправляет пользователям сообщения о создании, изменении и отмене бронирований.
\item Выполнено тестирование ключевых сценариев. Проверены бизнес-правила бронирования, расчет доступности, контрактные схемы REST и Kafka, а также клиентский сценарий от создания приглашения до отмены записи.
\end{enumerate}

Все основные задачи, поставленные во введении, выполнены. Разработанная система обеспечивает централизованное хранение данных о пользователях и записях, расчет доступных временных интервалов, защиту от конфликтующих бронирований, работу с пригласительными ссылками, интеграцию с внешним календарем и автоматическую отправку уведомлений.

Дальнейшее развитие платформы может включать расширение набора календарных провайдеров, развитие административной аналитики, настройку более гибких правил расписания, добавление платежных сценариев и расширение инструментов для групповой работы мастеров.
